\section{The homotopy types of the two Sweeney examples}\label{sec:monoids}

This section supplies maps realizing the homotopy equivalences.  In
particular, the identification does not infer a homotopy type from its
homology groups.  The argument treats Examples 5.2 and 5.3 of
\cite{Sweeney2025} at once.

\begin{definition}\label{def:Mr}
For an integer $r\geq 0$, let
\[
 M_r=\{1,x_{00},x_{01},x_{10},x_{11},y_1,\ldots,y_r\}.
\]
The identity is $1$.  Give $x_{ij}$ row $i$ and column $j$, and give each
$y_s$ row and column zero.  For two nonidentity elements $v,w$, set
\[
             vw=x_{\operatorname{row}(v),\operatorname{col}(w)}.
\]
Write $e=x_{00}$.  Thus the prompt's six-element monoid is $M_1$, with
$y_1=y$, and the monoid obtained by adjoining an identity to Sweeney's
Example 5.3 is $M_2$.
\end{definition}

\begin{theorem}[Explicit wedge decomposition]\label{thm:Mr-wedge}
Let $\mathcal W_r$ be the set of finite words, including the empty word,
on $\{1,\ldots,r\}$.  There are explicit based maps
$s_w:S^{|w|+2}\longrightarrow BM_r$ whose wedge is a homotopy equivalence
\[
 \Phi_r:\bigvee_{w\in\mathcal W_r}S^{|w|+2}\ \xrightarrow{\simeq}\ BM_r.
 \tag{A*}\label{eq:Mr-wedge}
\]
Consequently,
\[
 BM_1\simeq\bigvee_{n\geq2}S^n,
 \qquad
 BM_2\simeq\bigvee_{n\geq2}\bigvee_{1\leq j\leq2^{n-2}}S^n.
\]
The maps $s_w$ induce a primitive integral basis in every positive
homology degree in which homology is nonzero.
\end{theorem}

Here and below, the empty alphabet has just its empty word.  The proof
will use a short projective resolution to count the integral homology,
explicit cocycles to certify a basis, and explicit spheres realizing
that basis.

\subsection{The multiplication and an integral projective resolution}

The multiplication is associative: for three nonidentity elements both
parenthesizations have the row of the first and the column of the last.
Cases involving $1$ follow from its identity property.  Sending all
$y_s$ to $e$ and fixing the other elements defines a monoid retraction
$M_r\to M_0$.  The program
\texttt{code/A\_monoid\_verify.py} checks all multiplication triples for
$r=0,1,2$, as well as this retraction and all algebra identities used
below.

\begin{lemma}\label{lem:Mr-ring}
Put $R=\Z[M_r]$.  The following is an integral basis of $R$:
\[
 \begin{aligned}
 &e,\quad u=1-e,\quad a=x_{10}-e,\quad b=x_{01}-e,\\
 &c=x_{11}-x_{10}-x_{01}+e,\quad t_s=y_s-e\quad(1\leq s\leq r).
 \end{aligned}
\]
The elements $e,u$ are orthogonal idempotents with sum $1$.  The additive
subgroup $J$ spanned by $a,b,c,t_1,\ldots,t_r$ is an ideal.  Its only
nonzero product of basis elements is $ab=c$.  Its idempotent supports are
\[
 a=uae,\qquad b=ebu,\qquad c=ucu,\qquad t_s=ut_su.
\]
In particular $J^3=0$ and $R/J\cong\Z\times\Z$.
\end{lemma}
\begin{proof}
The inverse change of basis is integral:
$1=e+u$, $x_{10}=e+a$, $x_{01}=e+b$, $x_{11}=e+a+b+c$,
and $y_s=e+t_s$.  Expanding the products with the multiplication law in
Definition~\ref{def:Mr} gives the stated supports, $ab=c$, and zero for
all other products in $J$.  For example,
$(x_{01}-e)(x_{10}-e)=e-e-e+e=0$ and
$(y_s-e)(y_t-e)=e-e-e+e=0$.  These calculations also prove that $J$ is an
ideal and establish the quotient description.
\end{proof}

\begin{proposition}\label{prop:Mr-homology}
The augmented left $R$-module $T=\Z$ admits a projective resolution with
\[
 P_0=Re,\qquad P_1=Ru,\qquad
 P_n=(Re)^{r^{n-2}}\oplus(Ru)^{r^{n-1}}\quad(n\geq2).
\]
For $r=0$ this is interpreted as $P_2=Re$ and $P_n=0$ for $n>2$.
It follows integrally, in all degrees, that
\[
 H_0(BM_r;\Z)=\Z,\qquad H_1(BM_r;\Z)=0,\qquad
 H_n(BM_r;\Z)=\Z^{r^{n-2}}\quad(n\geq2).
\]
\end{proposition}
\begin{proof}
The augmentation sends $e$ to $1$ and $u,J$ to $0$.
Let $U$ denote the rank-one left module on which $u$ acts as $1$ and
$e,J$ act as zero.  Lemma~\ref{lem:Mr-ring} gives
\[
 Re=\Z e\oplus\Z a,\qquad
 Ru=\Z u\oplus\Z b\oplus\Z c\oplus\bigoplus_s\Z t_s.
\]
There are exact sequences
\[
 0\longrightarrow U\xrightarrow{\,1\mapsto a\,}Re
      \xrightarrow{\epsilon}T\longrightarrow0,
 \qquad
 0\longrightarrow Re\oplus U^{r}\longrightarrow Ru
      \xrightarrow{\,u\mapsto a\,}\Z a\longrightarrow0.
\]
For the second sequence, its kernel is
$\Z b\oplus\Z c\oplus\bigoplus_s\Z t_s$.
The map $Re\to Rb$ sending $e$ to $b$ sends $a$ to $c$, hence is an
isomorphism; each $\Z t_s$ is isomorphic to $U$.  This proves exactness
over $\Z$, including the asserted direct-sum decompositions of kernels.
Both $Re$ and $Ru$ are projective, being summands of $R$ defined by
idempotents.  At each step resolve every remaining $U$ summand by the
second exact sequence, and use the identity resolution on each $Re$
summand.  This gives the displayed $P_n$ and an exact projective
resolution in all degrees.  More concretely, from degree $2$ onward the
new $Re$ and $Ru$ summands map only into the preceding $Ru$ summands,
using $e\mapsto b$ and $u\mapsto t_s$.

Tensor on the left with the trivial right module $T$.  Then
$T\otimes_R Re=\Z$, while $T\otimes_R Ru=0$.  Every resulting
differential is zero: $d_1$ has zero source after tensoring, $d_2$ has
zero target, and all higher images lie in $Ru$ summands.  The claimed
homology follows from
$H_n(BM_r;\Z)=\operatorname{Tor}^{R}_n(T,T)$, equivalently the ordinary
bar resolution of the monoid.
\end{proof}

\begin{proposition}[Audit of the supplied periodic resolution]
\label{prop:Mr-Sweeney-resolution}
For $r=1$, Sweeney's displayed repeating resolution in Example 5.2 is
exact over $\Z$, including its periodic kernel identifications.
\end{proposition}
\begin{proof}
Use $y=y_1$.  Both left ideals $Rx_{00}$ and $Rx_{10}$ have the same
underlying submodule $Re$, with basis $e,x_{10}$.  The displayed maps,
before inserting zero in a first summand, are
\[
 \begin{aligned}
 D_1:R&\longrightarrow Re,& D_1(\rho)&=\rho(e-x_{10}),\\
 D_2:Re\oplus R&\longrightarrow R,&
 D_2(\xi,\rho)&=\xi x_{01}+\rho y,\\
 D_3:Re\oplus R&\longrightarrow R,&
 D_3(\xi,\rho)&=\xi(x_{10}-x_{11})+\rho(e-y).
 \end{aligned}
\]
Let $S\subset R$ be the free additive subgroup spanned by the five
nonidentity elements.  Direct multiplication gives
\[
 \ker D_1=S=\operatorname{im}D_2,
 \qquad \operatorname{im}D_3=\Z b\oplus\Z c\oplus\Z t_1.
\]
Indeed $D_1(1)=e-x_{10}$ and $D_1$ vanishes on $S$;
$\xi x_{01}$ spans $x_{01},x_{11}$ while $\rho y$ spans $y,e,x_{10}$.
Those five elements are distinct basis elements.  In $D_3$, its first
summand maps the basis $e,x_{10}$ to $-b,-b-c$, and its second summand
maps $1$ to $-t_1$ and every element of $S$ to zero.
The indicated independent integral bases show
\[
 \ker D_2=0\oplus\operatorname{im}D_3,
 \qquad \ker D_3=0\oplus\operatorname{im}D_2.
\]
For the first equality, the first summand of $D_2$ is forced to be
zero by its column-one coefficients; the kernel of right
multiplication by $y$ is precisely the displayed basis $b,c,t_1$.
For the second equality the independent $b,c,t_1$ coefficients force
the first summand and the coefficient of $1$ in the second to vanish.
Finally, the kernel of the augmentation $Re\to\Z$ is generated by
$x_{10}-e$, the image of $D_1$ up to sign.  This checks every exactness
assertion in the periodic resolution, with saturated integral images.
\end{proof}

\subsection{Explicit primitive cocycles and spherical cycles}

Use the normalized bar chains of $M_r$ with trivial coefficients.
Our bar convention multiplies consecutive entries in their written
order.  If the nerve convention takes the composite of consecutive
arrows $m,m'$ to be $m'm$, this bar simplicial set is $N(M_r^{\op})$.
Reversing the vertex order and barycentric coordinates gives the
canonical homeomorphism $B(M_r^{\op})\to BM_r$; every bar-space map
below is composed with this homeomorphism.  Thus the convention does
not change either the target homotopy type or the stated maps.
An
$n$-simplex is written $[m_1|\cdots|m_n]$, where all $m_i\ne1$, and
\[
 \begin{split}
 \partial[m_1|\cdots|m_n]={}&[m_2|\cdots|m_n]\\
 &+\sum_{i=1}^{n-1}(-1)^i
 [m_1|\cdots|m_im_{i+1}|\cdots|m_n]\\
 &+(-1)^n[m_1|\cdots|m_{n-1}].
 \end{split}
\]
Multilinear expressions in brackets denote the corresponding integral
linear combinations, not additional monoid elements.  Define functions
on the nonidentity elements by
\[
 \alpha(m)=\operatorname{col}(m),\quad
 \beta(m)=\operatorname{row}(m),\quad
 Y_s(m)=\begin{cases}1&m=y_s,\\0&\text{otherwise}.
 \end{cases}
\]
For a word $w=(s_1,\ldots,s_k)$ and $n=k+2$, put
\[
 \begin{split}
 f_w([m_1|\cdots|m_n])
   &=\alpha(m_1)\,Y_{s_1}(m_2)\cdots
                 Y_{s_k}(m_{n-1})\,\beta(m_n),\\
 z_w&=[(x_{01}-e)|(y_{s_1}-e)|\cdots|
                (y_{s_k}-e)|(x_{10}-e)].
 \end{split}
\]

\begin{lemma}\label{lem:Mr-duality}
Every $f_w$ is an integral cocycle, every $z_w$ is an integral cycle,
and, for words of the same length,
$f_w(z_v)=\delta_{wv}$.  The classes of the $z_w$ form an integral
homology basis; the classes of the $f_w$ form its dual cohomology basis.
\end{lemma}
\begin{proof}
For nonidentity $m,m'$, the product is a band element, and
\[
 \alpha(mm')=\alpha(m'),\quad
 \beta(mm')=\beta(m),\quad Y_s(mm')=0.
\]
In $\delta f_w$ the first two terms cancel by the first identity, and
the last two cancel by the second identity.  All other terms vanish by
the third identity.  This also covers the empty word, where there are
no remaining terms.

The two alternatives in each entry of $z_w$ have coefficient sum zero,
so both endpoint terms in its boundary vanish.  Every product of two
chosen adjacent actual elements is $e$, independently of both choices.
Hence every interior term in its boundary also vanishes.  Evaluation of
$f_w$ selects $x_{01}$ in the first entry, $x_{10}$ in the last, and the
specified $y_s$ at each interior entry; it gives the asserted Kronecker
pairing.

There are $r^{n-2}$ such cycles in degree $n$.  The pairing gives a split
injection from the free group on these cycles into $H_n$.  By
Proposition~\ref{prop:Mr-homology}, $H_n$ is free of exactly that rank,
so the split complement is zero.  The universal coefficient theorem,
with the preceding homology group free, gives the cohomology assertion.
\end{proof}

\begin{lemma}[The sphere map]\label{lem:Mr-spheres}
For each word $w$ there is a based map
$s_w:S^{|w|+2}\to BM_r$ taking an oriented fundamental class to $[z_w]$.
\end{lemma}
\begin{proof}
Put $n=|w|+2$.  Give $S^1$ a cell triangulation with two vertices, of
colors $0,1$, and two edges, both ordered from color $0$ to color $1$.
Form its join with $n-1$ copies of $S^0$, whose vertex colors are
$2,\ldots,n$.  Its realization is
$S^1*(S^0)^{*(n-1)}\cong S^n$.  There are $2^n$ top-dimensional
simplices, indexed by a choice of edge and one vertex in every $S^0$.
Index these choices by $\varepsilon=(\varepsilon_1,\ldots,
\varepsilon_n)\in\{0,1\}^n$.

Send the simplex indexed by $\varepsilon$ to the bar simplex whose
successive entries are selected by these bits from
\[
 (x_{01},e),\ (y_{s_1},e),\ldots,
                (y_{s_{n-2}},e),\ (x_{10},e).
\]
All vertices map to the unique vertex of the nerve.  The faces of
colors $0$ and $1$ pair the simplices obtained by flipping the first
bit.  Face $0$ deletes the first entry; face $1$ multiplies its first
two entries, always giving $e$.  Thus their images agree under the
pairing.  For $2\leq i<n$, the face of color $i$ pairs simplices by
flipping bit $i$; its bar image multiplies entries $i,i+1$, again
always giving $e$.  The last face deletes the last entry and respects
the last-bit pairing.  These checks prove that the simplex maps glue
to a continuous map on the whole sphere.  They also check all lower
face identifications, since those are composites of these face maps.

Orient the initial circle by the difference of its two edges and each
$S^0$ by the difference of its two vertices.  The resulting join
fundamental chain has coefficients $(-1)^{\sum_i\varepsilon_i}$.
Its image is precisely $z_w$.  Choosing any vertex as the base point
makes the map based.
\end{proof}

\begin{proof}[Proof of Theorem~\ref{thm:Mr-wedge}]
Wedge the maps of Lemma~\ref{lem:Mr-spheres}, using the CW weak topology.
This gives the map $\Phi_r$ in \eqref{eq:Mr-wedge}.
Lemma~\ref{lem:Mr-duality} says that it induces an integral homology
isomorphism.

The fundamental group of $BM_r$ has generators the nonidentity
elements and relations given by their products, as seen directly from
the $2$-skeleton of its nerve.  Every $x_{ij}$ is idempotent and hence
maps to $1$ in that group.  The relation $y_se=e$ then sends every
$y_s$ to $1$ as well.  Thus $BM_r$ is simply connected.  The source of
$\Phi_r$ is a wedge of spheres of dimension at least two and is also
simply connected.  The homological Whitehead theorem now proves that
$\Phi_r$ is a homotopy equivalence.
\end{proof}

This explicitly settles the first possibly problematic case in A2:
for $r\geq1$, choose any word of length one and a based homotopy inverse
$\Psi_r$ of $\Phi_r$.  The composite
\[
 BM_r\xrightarrow{\Psi_r}\bigvee_{w}S^{|w|+2}
         \xrightarrow{\text{projection}}S^3
\]
has degree $1$ on the corresponding $H_3$ generator.  The same
construction gives a retraction onto every displayed sphere.

\subsection{Products, operations, and the auxiliary assertions of R3}

\begin{proposition}\label{prop:Mr-cohomology}
For $A=\Z$ or $\F_p$, $H^0(BM_r;A)=A$ and the positive-degree
cohomology has basis $f_w$ in degree $|w|+2$.  Every product of two
positive-degree classes is zero.  Every Steenrod square of positive
degree is zero; at odd primes every positive-degree reduced power and
the Bockstein are zero.  Every ordinary triple Massey product of
positive-degree classes is defined and is the singleton $\{0\}$.
\end{proposition}
\begin{proof}
The additive assertion follows from Proposition~\ref{prop:Mr-homology}
and the universal coefficient theorem.  The wedge decomposition
already proves the product assertion by restricting to each sphere.
Here is also a direct integral bar-cochain check, which in particular
checks the products requested in R3 before reduction modulo a prime.
For words $v,w$ of lengths $k,l$, put $p=k+2$, $q=l+2$, and let
$\gamma$ be the indicator function of $x_{11}$.  Define a cochain of
degree $p+q-1$ by
\[
 \begin{split}
 h_{v,w}([m_1|\cdots|m_{p+q-1}])={}&
 \alpha(m_1)\prod_{i=1}^{k}Y_{v_i}(m_{i+1})\,
 \gamma(m_p)\\
 &\hspace{-10mm}\cdot\prod_{j=1}^{l}Y_{w_j}(m_{p+j})\,
 \beta(m_{p+q-1}).
 \end{split}
\]
For nonidentity $a',b'$, one has
$\gamma(a'b')=\beta(a')\alpha(b')$.
All terms in the coboundary of $h_{v,w}$ cancel in endpoint pairs or
vanish at a $Y$ slot, except the term merging the two central entries.
Consequently
\[
       \delta h_{v,w}=(-1)^p f_v\smile f_w.
\]
This gives explicit null cochains for every positive product, integrally
and after coefficient reduction.

For a positive-degree Steenrod operation on a class of degree $d$, its
output degree is larger than $d$.  On a sphere summand of that output
dimension the input restricts to zero, so naturality makes the output
zero.  Those restrictions jointly detect cohomology, and hence the
operation is zero.  This proves all the stated squares and powers;
the same argument proves the Bockstein assertion (which also follows
from torsion-free integral homology).

All triple Massey products of positive-degree classes are defined
because the required products are zero.  Their indeterminacy is a sum
of positive-degree products, hence zero.  If the input degrees are
$d_1,d_2,d_3$, their output degree $D=d_1+d_2+d_3-1$ is larger than each
input degree.  Restrict to any $S^D$ summand.  All three inputs restrict
to zero and their Massey product there is $\{0\}$.  Naturality and the
jointly injective restrictions in degree $D$ prove the result.  In
particular the first possible product, the degree-five product
$\langle f_{\varnothing},f_{\varnothing},f_{\varnothing}\rangle$, is
exactly $\{0\}$.  Alternatively the displayed primitives give a
defining system whose Massey cocycle evaluates to zero on every
$z_w$: each summand contains a factor forcing an interior entry to be
$x_{11}$, whereas those cycles have only $e,y_s$ at interior entries.
\end{proof}

\begin{proposition}\label{prop:Mr-auxiliary}
The monoid $M_r$ has trivial universal group.  The inclusion
$M_0\hookrightarrow M_r$ and the retraction $y_s\mapsto e$ exhibit
$BM_0\simeq S^2$ as a homotopy retract.  For $r=1$, the rational algebra
$\Q[M_1]$ has exactly two simple modules, its trivial simple has first
syzygy the other simple in the projective cover $\Q[M_1]e$, and
\[
 \Q[M_1]/\Q\{x_{00},x_{01},x_{10},x_{11}\}
       \cong\Q[t]/(t^2).
\]
Every phantom map from $BM_r$ to a CW target is null-homotopic.
\end{proposition}
\begin{proof}
The universal-group and retraction assertions were proved above.
Over $\Q$ the ideal $J$ of Lemma~\ref{lem:Mr-ring} is nilpotent and the
quotient is the semisimple algebra $\Q\times\Q$.  Thus $J$ is the
Jacobson radical and the two idempotent vertices give exactly the two
simple modules.  The first exact sequence in
Proposition~\ref{prop:Mr-homology} computes the stated syzygy.  The
band ideal has basis $e,a,b,c$; its quotient has basis $u,t$ with $u$
the identity and $t^2=0$.

For the phantom assertion, compose a phantom map with $\Phi_r$.
Every sphere restriction is null, since its map to $BM_r$ factors
through a finite subcomplex: a map from a finite CW complex into a CW
complex is homotopic to one whose image lies in a finite subcomplex.
A based map out of a CW wedge is determined up to based homotopy by
its restrictions to the summands.  Hence the composite is null and
so is the original map.  For an unbased phantom map, choose its
basepoint value and use that the source is simply connected; the
usual $\pi_1$-action cannot turn a nonzero based class into the zero
class.
\end{proof}

\subsection{A finite two-object structural model}

\begin{proposition}\label{prop:Mr-Karoubi}
Let $\mathcal K_r$ be the full subcategory of $\Kar(M_r)$ on $1,e$.
Its morphism sets are
\[
 \begin{array}{c|cc}
  \Hom(\text{row},\text{column})&1&e\\\hline
  1&M_r&\{e,x_{01}\}\\
  e&\{e,x_{10}\}&\{e\}.
 \end{array}
\]
Composition is multiplication in $M_r$, with $e$ the identity of the
second object.  The inclusion of the object $1$ induces a homotopy
equivalence $BM_r\to B\mathcal K_r$.
\end{proposition}
\begin{proof}
The formula $\Hom(d,d')=d'M_rd$ gives the table and all composition
laws.  Apply Quillen's Theorem A, in its dual formulation
\cite[p.~93]{Quillen1973}, to the inclusion $i:M_r\to\mathcal K_r$.
The comma category $(i\downarrow1)$ has terminal object $\id_1$.
An object of $(i\downarrow e)$ is an element $f\in eM_r$; a morphism
$f\to g$ is an $m\in M_r$ with $gm=f$.  Let $F$ be the constant
functor with value the object $e$ of this comma category.  For each
$f$, the element $f\in M_r$ is a morphism $f\to e$, since $ef=f$.
These components form a natural transformation from the identity to
$F$: for $m:f\to g$, its naturality equation is exactly $gm=f$.
Thus the comma category is contractible.  Both comma nerves are
contractible, proving the assertion.
\end{proof}

This supplies the finite structural model requested in A3, together
with its comparison map and proof.  No use of the recent
$M(Y,f)$ construction is needed.  The same category and the same proof
work for the second Sweeney example.  Its different homotopy type is
accounted for by the second symbol in each word, not just by an
observed finite list of Betti numbers.

\subsection{Independent computational checks and scope}

The attached verifier uses only Python's standard library.  For
$r=0,1,2$ it checks the full multiplication table, the integral change
of basis and every product in Lemma~\ref{lem:Mr-ring}, the module
identities giving the projective-resolution kernels, and the monoid
retraction.  Through degree eight it checks every simplex-face
identification in every sphere map and every boundary $\partial z_w$;
it checks the integral dual pairings and the cup-product pairings.
It also exhaustively verifies the cocycle equation on all normalized
simplices through cochain degree four and the explicit integral
cup primitives for bidegrees $(2,2),(2,3),(3,2)$.
The all-degree proofs above,
not extrapolation from these checks, establish the homotopy types.
